% Original PDF page 22; hash in recovery-map.json.
\recoverypage{22}
\section*{N  Retrieval substrates and premise formation}
\noindent
The architecture has more than one retrieval implementation. This appendix separates lexical lookup,\linebreak[4]
embedding-neighbor search, graph navigation, and formal premise admission. It does not claim that\linebreak[4]
a single request automatically combines every index or that a learned fusion policy has been trained.\linebreak[4]
The relevant inspected sources are the legal sample builder, GraphRAG thin client, FAISS store, and\linebreak[4]
Hammer selector contract. [artifact S22,S34--S36]\par

\subsection*{N.1  BM25 in sample construction and corpus lookup}
\noindent
The sample builder invokes a BM25 frame selector to produce candidates and a selected ontology\linebreak[4]
frame. This is supervision provenance: a selected frame may be useful but is not automatically\linebreak[4]
a human-validated semantic label. The GraphRAG thin client is a different consumer of lexical\linebreak[4]
indexing. It tokenizes the query, locates relevant posting shards, and computes a term contribution of\linebreak[4]
the form\par

% Newly transcribed from original PDF equation (15); see recovery-map.json.
\begin{equation}
s(t,d)=\operatorname{idf}(t)\frac{\operatorname{tf}(t,d)(k_1+1)}{\operatorname{tf}(t,d)+k_1(1-b+b\,|d|/\overline{|d|})}.\tag{15}\label{eq:original-15}
\end{equation}

\noindent
For one inspected posting schema, term frequency combines configured title and body weights.\linebreak[4]
However, a separate schema branch keyed by legal identifiers sums the supplied term frequencies\linebreak[4]
instead. Consequently, a response labeled ``bm25'' is not by itself evidence that every data schema\linebreak[4]
used the same full scoring formula. The manifest, posting schema, scoring branch, tokenizer, and\linebreak[4]
parameters belong in a reproducible run record. [artifact S34]\linebreak[4]
Lexical retrieval is useful for citations, names, identifiers, and distinctive vocabulary. It can also\linebreak[4]
miss paraphrases, rank a superficially matching exception above a controlling rule, or omit a relevant\linebreak[4]
provision. Its score supplies ordering, not entailment, completeness, or applicability. Frame-selection\linebreak[4]
errors can propagate into compiler-derived supervision unless independent checks expose them.\par

\subsection*{N.2  Vector indexing is separate from representation learning}
\noindent
The FAISS store maintains numerical indexes, per-collection metadata, and mappings between\linebreak[4]
external identifiers and internal index positions. Its inspected constructor requires real FAISS and\linebreak[4]
NumPy availability; legacy mock definitions in the file are not evidence of a successful real store\linebreak[4]
query. An evaluated route must identify the actual backend and status. [artifact S35]\linebreak[4]
Three operations must be recorded separately: producing a vector with a frozen or trained encoder;\linebreak[4]
fitting any index structures that a selected index type requires; and changing semantic model parame-\linebreak[4]
ters. The adaptive autoencoder's embedding reconstruction is a fourth, separately specified numerical\linebreak[4]
operation. A published numerical vector does not prove that the corresponding source was formally\linebreak[4]
interpreted, and a store over vectors does not establish that those vectors came from this autoencoder.\linebreak[4]
The thin client declares a vector subcommand and contains an embedding helper, but its inspected\linebreak[4]
dispatch does not execute a vector query. Actual vector search must therefore be attributed to another\linebreak[4]
verified entry point rather than to that declaration. This is a route-specific gap, not a claim that the\linebreak[4]
repository lacks vector indexing. [artifact S34,S35]\par

\subsection*{N.3  Graphs connect evidence without silently adding axioms}
\noindent
The thin client can retrieve incoming/outgoing adjacency records and execute bounded walks over\linebreak[4]
CID-linked nodes. Its results retain neighbor identifiers and edge metadata. A graph can connect\linebreak[4]
a source passage to a definition, entity, related artifact, or dependency; an index edge does not\linebreak[4]
necessarily assert a logical implication. [artifact S34]\linebreak[4]
At the query boundary, lexical and vector hits can nominate seeds, graph navigation can expand their\linebreak[4]
context, and exact artifact identities can retrieve declarations. This is a proposed composition of the\linebreak[4]
available interfaces, not a measured fusion algorithm. A premise-admission stage must then inspect\linebreak[4]
source lineage, entity and temporal alignment, domain applicability, declared authority, and the\linebreak[4]
relevant interpretation. Keep ``retrieved'', ``admitted premise'', ``assumed'', and ``proved'' as separate\linebreak[4]
states.\par

